\section{Conclusion}
\label{sec:conclusion}

This study considered four questions separately: whether a configuration is identifiable, whether an implementation completed, whether it carries predictive signal, and whether its gain meets the rule for moving to the next stage.  In the UrbanEV registry, 865 of 1,224 configurations have completed artifacts, but only 156 meet S1.  Completion is useful information, but it is not the same as exact reproduction.

The later model experiments produced different results at each stage.  The target-dependent oracle reduced RMSE by 16.926\% relative to TimeXer.  Strict prequential fixed fusion retained a 3.178\% internal gain and increased batch-one P50 latency by 48.86\% on the frozen device.  Learned routing, the Chronos-2 teacher screen, aligned residual distillation, and the Paris development teacher were each evaluated under their registered rules, but none met its progression rule.  Positive router and Paris intervals did not replace the project-specific 1\% point requirement.

Metric locks, full control reruns, prediction-identity checks, and time-stamped stopping rules made each negative result traceable to a prespecified comparison.  Paris formal and protected targets were kept out of the present analysis.  The separate store contains 48 non-self-referential files, including 29 recoverable Git files, and the protected analytical-access count remains zero.  This is a documented storage rule; it is not encrypted storage or a protected-test result.

The main result concerns evaluation practice rather than a new model architecture.  In this case, separate checks make clear which conclusions the stored evidence supports.  The wording rules and audit checklist have not yet been tested by another team or on another benchmark.  Future work should apply a newly frozen model and loss specification to an untouched data boundary, use multiple seeds and hardware environments, and test the same stopping rules independently.  At present, the evidence supports an internal accuracy--cost result and the recorded negative decisions, while formal and protected targets remain unused for later evaluation.  It does not support a state-of-the-art or deployment-readiness claim.
